\documentclass[10pt,a4paper]{amsart}
\usepackage[margin=19mm]{geometry}
\usepackage{amsmath,amssymb,mathtools}
\usepackage[scr=rsfs,cal=euler]{mathalfa}
\usepackage{xfrac}
\usepackage[unicode,hidelinks,bookmarks=false]{hyperref}
\newcommand{\ie}{i.e.\ }
\newcommand{\cf}{cf.\ }
\newcommand{\resp}{respectively\ }
\newcommand{\Subsection}{\subsection}
\providecommand{\nobreakdash}{\nobreak\hbox{-}}
\setlength{\emergencystretch}{2em}
\multlinegap0pt
\usepackage{xcolor}
\usepackage{dsfont}
\usepackage{amssymb}
\usepackage{enumerate}
\usepackage{mathtools}
\newtheorem{thm}{Theorem}
\newtheorem{cor}{Corollary}
\title{$\beta$-Uniform Convexity and Divisible Domains}
\author{Amelia Pompilio}
\date{Source: arXiv:2410.20071v2 (CC BY 4.0)}
\begin{document}
\maketitle

\noindent\emph{Source and licence.} $\beta$-Uniform Convexity and Divisible Domains. Version arXiv:2410.20071v2 (CC BY 4.0), distributed by arXiv under the Creative Commons Attribution 4.0 International licence, http://creativecommons.org/licenses/by/4.0/, as recorded in the arXiv licence metadata for this exact version and shown on the abstract page https://arxiv.org/abs/2410.20071v2. Full licence notice: \texttt{licenses/arxiv-2410.20071-CC-BY-4.0.txt}.\par\smallskip

\noindent\textit{Editorial scope.} Counted unit: the article's cor cor:mainfinsler (source0.tex lines 877--879). The article's complete original proof of it is delivered with it (source0.tex lines 881--1013, one \texttt{proof} environment of 133 lines). No mathematics is written, repaired or re-derived: the statement and the proof are the article's own lines, in the article's own notation, and no retained line is substituted or re-flowed.  Retained uncounted as the source-local context the proof needs: lines 756--759; lines 865--867.  Nothing was omitted from any retained span, so the package carries no omission record.  No retained line contains a citation, so the wrapper carries no bibliography and no bibliography entry.  No retained line contains a figure environment, an image-inclusion command or drawing code, so no image is delivered and the package declares no asset dependency.  Editorial formatting only, in addition: an AMS \texttt{amsart} preamble that restates the article's own declarations and macros used by the retained lines, namely the environments \texttt{thm}, \texttt{cor} on separate counters, together with the standard packages those lines need.  The retained environments print in this document as Theorem 1, Corollary 1 in order of appearance, whereas the article prints the same items with the numbering of its own full document.  The complete native source of the article as distributed in the arXiv e-print for this exact version is bundled unchanged as \texttt{sources/167443/source0.tex}.
\begin{thm}
\label{thm:mainminkowski}
Let $\Omega \subset \mathbb{RP}^n$ be a strictly convex divisible domain. Then for any $x_0 \in \Omega,$ the space $\big(T_{x_0}\Omega, M_\Omega(x_0, \cdot)\big)$ is $\beta$-uniformly convex for some $\beta \in [2, \infty)$.
\end{thm}

\begin{equation}\label{eqn:betauc}
    \delta_\Omega(\varepsilon) \geq C\varepsilon^\beta
\end{equation}

\begin{cor}\label{cor:mainfinsler}
    Let $\Omega \subset \mathbb{RP}^n$ be a strictly convex divisible domain. Then the space $\big(T_{x_0}\Omega, F_\Omega(x_0, \cdot)\big)$ is $\beta$-uniformly convex for some $\beta \in [2, \infty)$.
\end{cor}

\begin{proof}[Proof of Corollary \ref{cor:mainfinsler}]

 We will show that 
 \begin{equation}\label{FisBUC}
        \delta(\varepsilon) = \inf \bigg(1-\frac{F_\Omega(x_0,x+y)}{2}\bigg) \geq C_0 \varepsilon^\beta
    \end{equation}
    for $\varepsilon \in [0,2],$ where the infimum is taken over all $x,y \in T_{x_0}\Omega$ with $F_\Omega(x_0,x), F_\Omega(x_0,y) = 1$, and $F(x_0,x-y) \geq \varepsilon$.

We use the notation 
\begin{equation*}
     M_\Omega^{\pm}(v) \coloneq M_\Omega(\pm v),  v \in T_x\Omega
 \end{equation*}
 throughout.
 There exists an $A > 0$ such that
\begin{equation}\label{equivalencefwdbkwd}
   \frac{1}{A}M_\Omega^-(x_0,\cdot) \leq M_\Omega^+(x_0,\cdot) \leq AM_\Omega^-(x_0,\cdot).
 \end{equation}
 for all $v \in T_x\Omega$. Recall that the Hilbert-Finsler norm $F_\Omega(x_0, \cdot)$ is the symmetrization of the forward and backward Minkowski functionals, that is
 \begin{equation}\label{FisSymmetrization}
     F_\Omega(x_0,\cdot) = \frac{M^+_\Omega(x_0,\cdot) + M^-_\Omega(x_0,\cdot)}{2}.
 \end{equation}
 Using \ref{equivalencefwdbkwd}, we have
 \begin{equation}
  F_\Omega(x_0,\cdot) \leq \frac{M^+_\Omega(x_0,\cdot)}{2} +    \frac{AM^+_\Omega(x_0,\cdot)}{2} = \frac{A+1}{2}M^+_\Omega(x_0,\cdot),
 \end{equation}
and since $F_\Omega(x_0,x-y) \geq \varepsilon$, this implies that
 \begin{equation}\label{unnormalizedlower}
  M^+_\Omega(x_0,x-y) \geq \frac{2\varepsilon}{A+1}.   
 \end{equation}
 And similarly, since $F_\Omega(x_0,x)=F_\Omega(x_0,y)=1,$ we have
 \begin{equation}\label{lowerboundforForward}
   M^+_\Omega(x_0,x) \geq \frac{2}{A+1}  \text{ and } M^+_\Omega(x_0,y) \geq \frac{2}{A+1}.
 \end{equation}
We wish to use the lower bound on the modulus of convexity that we found in \ref{eqn:betauc} for the Minkowski functional in order to find the desired lower bound in \ref{FisBUC} for the Hilbert-Finsler norm. The lower bound in \ref{eqn:betauc} holds for $\varepsilon$-separated vectors $x,y \in T_{x_0}\Omega$ whose forward Minkowski functional values are both less than or equal to 1. Since we are assuming that $F_\Omega(x_0,x)=F_\Omega(x_0,y)=1$, it follows from \ref{FisSymmetrization} that $M^+_\Omega(x_0,\cdot)$ or $M^-_\Omega(x_0,\cdot)$ may be greater than 1. Thus in order to use the lower bound in \ref{eqn:betauc}, we first normalize $x$ and $y$ with respect to $M^+_\Omega(x_0,x)$ and $M^+_\Omega(x_0,y)$.

So let $a = M^+_\Omega(x_0,x)$ and $b = M^+_\Omega(x_0,y)$, and let $a \geq b$. Then $M^-_\Omega(x_0,x) = 2-a$ and $M^-_\Omega(x_0,y) = 2-b$. Define
\begin{equation}
    \overline{x} := \frac{x}{M^+_\Omega(x_0,x)}=\frac{x}{a} \text{ and } \overline{y} := \frac{y}{M^+_\Omega(x_0,y)}=\frac{y}{b}.
\end{equation}
Then $M^+_\Omega(x_0, \overline{x}) = M^+_\Omega(x_0, \overline{y}) = 1$.

Next, we aim to show that $M^+_\Omega(x_0, \overline{x}-\overline{y})$ is bounded below by some multiple of $\varepsilon$. This will allow us to apply Theorem \ref{thm:mainminkowski} to $M^+_\Omega(x_0,\frac{\overline{x}+\overline{y}}{2})$.

Subtracting the backward Minkowski functionals of the normalized $x$ and $y$, we get
\begin{equation}
 M^-_\Omega(x_0, \overline{y})-M^-_\Omega(x_0, \overline{x})= \frac{2-b}{b} - \frac{2-a}{a} = \frac{2(a-b)}{ab}.  
\end{equation}
By the triangle inequality, we have
\begin{equation}
 M^-_\Omega(x_0, \overline{y})-M^-_\Omega(x_0, \overline{x}) \leq M^-_\Omega(x_0, \overline{y}-\overline{x}) = M^+_\Omega(x_0, \overline{x}-\overline{y})
\end{equation}
Then combining the previous two lines gives
\begin{equation}\label{a-b}
    a-b \leq \frac{ab}{2} M^+_\Omega(\overline{x}-\overline{y}) \leq 2 M^+_\Omega(\overline{x}-\overline{y}),
\end{equation}
where the last estimate follows since $a,b \leq 2.$
On the other hand,
\begin{equation}
    x-y = a\overline{x} - b\overline{y} = a\overline{x}-b\overline{x}+b\overline{x}-b\overline{y} = (a-b)\overline{x}+b(\overline{x}-\overline{y}).
\end{equation}
Applying $M^+_\Omega(x_0,\cdot)$, we get
\begin{equation}
    M^+_\Omega(x_0,x-y) \leq (a-b)M^+_\Omega(x_0,\overline{x}) + bM^+_\Omega(x_0,\overline{x}-\overline{y})
\end{equation}
and since $M^+_\Omega(x_0,\overline{x})=1$ and $b \leq 2,$ using the estimate in \ref{a-b} gives
%\begin{equation}
\begin{align}
 M^+_\Omega(x_0,x-y) & \leq bM^+_\Omega(x_0,\overline{x}-\overline{y}) + (a-b)\\
 & \leq 2M^+_\Omega(x_0,\overline{x}-\overline{y}) + 2M^+_\Omega(x_0,\overline{x}-\overline{y})\\
 & = 4M^+_\Omega(x_0,\overline{x}-\overline{y}).
 \end{align}
So by \ref{unnormalizedlower} and the previous line, we get
\begin{equation}
    M^+_\Omega(x_0,\overline{x}-\overline{y}) \geq \frac{\varepsilon}{2(A+1)}=\varepsilon_0
\end{equation}
and hence, by Theorem \ref{thm:mainminkowski},
\begin{equation}
    1-M^+_\Omega\bigg(x_0,\frac{\overline{x}+\overline{y}}{2}\bigg) \geq C\varepsilon_0^\beta.
\end{equation}

We now find bounds for $M^+_\Omega(x_0,\frac{x+y}{2})$ and $M^-_\Omega(x_0,\frac{x+y}{2})$. Taking their average as in \ref{FisSymmetrization} will give the desired result. We begin by observing that
\begin{equation}
    \frac{x+y}{2}=\frac{(a-b)\overline{x}}{2} + \frac{b(\overline{x}+\overline{y})}{2}.
\end{equation}
Applying $M_\Omega^+(x_0,\cdot)$ and simplifying gives
\begin{equation}
 M^+_\Omega\bigg(x_0,\frac{x+y}{2}\bigg)  \leq \frac{a+b}{2} - Cb\varepsilon_0^\beta .
\end{equation}
For the reverse Minkowski functional, ordinary convexity gives
\begin{align}
    M^-_\Omega(x_0,\frac{x+y}{2}) &\leq \frac{M_\Omega^-(x_0,x) + M_\Omega^-(x_0,y)}{2}\\
    & = \frac{(2-a) + (2-b)}{2}\\
    & = 2 - \frac{(a+b)}{2}.
\end{align}
Finally, we take the average
\begin{align}
 F_\Omega\bigg(x_0,\frac{x+y}{2}\bigg) & = \frac{M^+_\Omega(x_0,\frac{x+y}{2}) + M^-_\Omega(x_0,\frac{x+y}{2})}{2} \\
 & \leq \frac{\big(\frac{a+b}{2} - Cb\varepsilon_0^\beta\big) + \big(2 - \frac{a+b}{2}\big)}{2} \\
 & = 1-\frac{Cb\varepsilon_0^\beta}{2}.
\end{align}
And since \ref{lowerboundforForward} implies that $b \geq \frac{2}{A+1}$, we have
\begin{equation}
    1-F_\Omega\bigg(x_0,\frac{x+y}{2}\bigg) \geq \frac{C}{2^\beta(A+1)^{\beta+1}}\varepsilon^\beta
\end{equation}
as desired.
%\end{equation}
% We use the notation
% \begin{equation*}
%     M_\Omega^{\pm}(v) \coloneq M_\Omega(\pm v)
% \end{equation*}
% throughout. There exists an $A > 0$ such that
% \begin{equation}
%     \frac{1}{A}M_\Omega^- \leq M_\Omega^+ \leq AM_\Omega^-.
% \end{equation}

% Take $x,y \in \Omega$ such that $F_\Omega(x-y) > \varepsilon$. Since

% \begin{equation*}
%     F_\Omega(x,v) = \frac{M^+_\Omega(x,v) + M^-_\Omega(x,v)}{2},
% \end{equation*}
% without loss of generality we have $M^+_\Omega(x-y) > \varepsilon$, and so $M^-_\Omega(x-y) > \frac{\varepsilon}{A}$. Then by \ref{eqn:betauc},
% \begin{equation}
%     1 - \frac{M^+_\Omega(x+y)}{2} \geq C_+\varepsilon^\beta
% \end{equation}
% and
% \begin{equation}
%     1 - \frac{M^-_\Omega(x+y)}{2} \geq C_-\bigg(\frac{\varepsilon}{A}\bigg)^\beta.
% \end{equation}
% Then
% \begin{equation}
%     1-\frac{F_\Omega(x+y)}{2} = 1-\frac{\frac{M^+_\Omega(x+y) + M^-_\Omega(x+y)}{2}}{2} > \Bigg(\frac{C_+ + \frac{C_-}{A^\beta}}{2}\Bigg)\cdot \varepsilon^\beta > 0
% \end{equation}
\end{proof}

\end{document}
